\documentclass[10pt]{article}
\usepackage[margin=20mm]{geometry}
\usepackage{amsmath,amssymb,mathtools,fontspec,unicode-math,xurl,hyperref}
\IfFontExistsTF{DejaVu Serif}{\setmainfont{DejaVu Serif}}{\setmainfont{Latin Modern Roman}}
\setmathfont{latinmodern-math.otf}
\setlength{\parindent}{0pt}\setlength{\parskip}{6pt}\sloppy
\title{Local obstruction data and signature conditions in K3 realization}
\author{AI-assisted technical review}
\date{}
\begin{document}
\maketitle
Abstract: The supplied manuscript contains a failure of choice independence in its local obstruction construction. An explicit example over \(\mathbb Q_3\), satisfying the stated global numerical conditions, gives two admissible decompositions with different obstruction values. A second explicit counterexample disproves a fixed-characteristic-polynomial realization statement in degree twenty. Four further findings concern a \(2\)-adic determinant quantifier, the abstract characterization of signature maps, dimensions in two constructions, and an omitted root hypothesis. These findings are distinguished from the geometric realization and finiteness arguments that remain compatible with the checks below. The assessment is Major; the scientific decision remains pending.\par
\section*{Source identity and scope}
The reviewed source is Eva Bayer-Fluckiger, Automorphisms of K3 surfaces, signatures, and isometries of lattices, arXiv:2209.06698v4, posted 1 May 2024, with internal date 2 May 2024, 59 pages [main]. The journal identity is DOI 10.4171/JEMS/1598 (https://doi.org/10.4171/JEMS/1598), online first 30 January 2025 [publisher]. The author manuscript has not been established to be identical to the publisher version. All page references below are one-based PDF pages, equal to the printed page numbers.\par
The review reads the full manuscript, including definitions, local quadratic and hermitian forms, the obstruction argument, lattice surgery, K3 realization, examples and finiteness. It supplies exact counterexamples to specified statements, rather than a counterexample to every main result. Source images used for the findings and selected positive checks were actually viewed; their precise page list and hashes are in the accompanying provenance. Imported classification, Torelli and finiteness theorems are identified as dependencies, not represented as completely reproved here.\par
\section*{The local obstruction is not independent of choices}
Problem 1 (Major).  Notation 7.2 (p. 19) misses a possible local exchange. Consequently Proposition 8.1 and Lemma 8.3 (pp. 21–22) fail, and the ensuing definition of \(\epsilon^p\) is not independent of the chosen local data. This affects the obstruction used in Theorems 9.7 and 12.1 (pp. 25, 29).\par
\subsection*{An input satisfying the numerical hypotheses}
Set \[f(X)=X^2-7X+1,\qquad F(X)=f(X)^2(X+1)^4.\] Give the \(f\)-part signature \((2,2)\) and the \((-1)\)-part signature \((2,2)\), so the total signature is \((4,4)\). Then \[F(0)=1,\quad \deg F=8,\quad F(1)=400,\quad F(-1)=0.\] Thus both (C0) and (C1) hold. The determinant prescribed in Section 11 for the \((-1)\)-part is \[D_-=|f(-1)^2|=81=1\quad\text{in }\mathbb Q^\times/\mathbb Q^{\times2}.\] The field \(E_f=\mathbb Q(\sqrt5)\) is ramified only at \(5\), whereas \(f\bmod5=(X-1)^2\). Notation 7.2 therefore gives \[\Pi_{f,X+1}=\varnothing.\] For \(I=\{f,X+1\}\), its group \(C_0(I)\) permits the function \(c(f)=1,\ c(X+1)=0\).\par
\subsection*{Two decompositions of the same local quadratic space}
At \(p=3\), the extension \(E_f/\mathbb Q_3\) is unramified quadratic. Multiplication by \[\alpha=(7+3\sqrt5)/2\] reduces to \(-1\) on the natural integral basis. The trace forms of Notation 1.2 are \[b_\lambda(x,y)=\operatorname{Tr}_{\mathbb Q_3(\sqrt5)/\mathbb Q_3}
       \bigl(\lambda x\sigma(y)\bigr)
       \simeq \langle2\lambda,-10\lambda\rangle.\] Consider \[\begin{array}{c|cc}
 & f\text{-part}&(-1)\text{-part}\\ \hline
 A&b_1\perp b_1&\langle1,1,1,1\rangle\\
 B&b_3\perp b_1&\langle3,3,1,1\rangle .
\end{array}\] On the \(f\)-part, let the generator act by multiplication by \(\alpha\) on each copy; on the other part let it act by \(-1\). Both characteristic polynomials are \(F\).\par
All four blocks have determinant squareclass \(1\). At \(3\), the Hasse invariant of each block in row \(A\) is zero. Row \(B\) changes each block’s Hasse invariant by \(1\): indeed \[(6,-30)_3=-1,\qquad (3,3)_3=-1,\] while the unit-coefficient contributions and the cross terms do not change the comparison. Therefore both total quadratic spaces have dimension eight, determinant \(1\) and Hasse invariant zero: both are \(H^4\) over \(\mathbb Q_3\).\par
The equivariant residue of row \(A\) is zero. In row \(B\) the nonzero residues come from the coefficients with odd \(3\)-adic valuation: \[\partial(b_3\perp b_1)=
      \langle-1,-1\rangle_{\chi_-},\qquad
 \partial\langle3,3,1,1\rangle=
      \langle1,1\rangle_{\chi_-}.\] They lie in the same \(\chi_-\) summand, because \(\alpha\equiv-1\bmod3\), and cancel. In particular both rows belong to the same set \(C^3_0\) used for \(\delta=0\). The class \(\langle1,1\rangle\) over \(\mathbb F_3\) is nonzero: \(x^2+y^2=0\) has no nonzero solution.\par
Relative to the reference forms of Sections 5–6, the two rows give \[a=(0,0),\qquad a'=(1,1)\in C^3_0.\] The function \(c\) above consequently satisfies \[\epsilon_a(c)=0,\qquad \epsilon_{a'}(c)=1.\] This is a direct contradiction to Lemma 8.3. It is not a failure confined to an arbitrary nonzero residue, an inadmissible signature, or an absent multiplicity: the example has \(\delta=0\), satisfies (C0)–(C1), and uses the allowed multiplicity \(n_f=2\).\par
\subsection*{The reverse defect and the extent of repair}
The stated edge condition can also add an impossible exchange. For a compact corroborating example, put \[g(Y)=Y^2+3Y+8,\quad
 h(X)=X^2g(X+X^{-1})=X^4+3X^3+10X^2+3X+1.\] The discriminant of \(g\) is \(-23\) and the norm of \(r^2-4\), for a root \(r\) of \(g\), is \(g(2)g(-2)=108\), not a rational square; hence \(h\) is irreducible. Over \(\mathbb Q_3\) the roots \(r_+,r_-\) of \(g\) near \(2,-2\) satisfy \[v_3(r_+-2)=2,\quad (r_+-2)/9\equiv1\pmod3,\qquad
 v_3(r_-+2)=1.\] Thus \(r_+^2-4\) is a square, while \(r_-^2-4\) has squareclass \(3\). The algebra is \(\mathbb Q_3\times\mathbb Q_3\times\mathbb Q_3(\sqrt3)\); its only ramified quadratic block has reduction \(\alpha=-1\). Yet \(h\bmod3=(X-1)^2(X+1)^2\). For \(h(X)(X-1)^4\), Notation 7.2 adds the edge to \(X-1\), although a Hasse change in the nonsplit block changes only the \(\chi_-\) residue. The \(+1\) block cannot cancel it. This also disproves the stability assertion in Proposition 8.1.\par
A repair must classify actual local residue exchanges, including nonsplit unramified blocks whose chosen generator reduces to \(\pm1\), and then reprove the orbit and choice-independence assertions. Merely requiring a ramified place with the matching sign removes the second example but does not repair the first. The obstruction group and its global use must be reconsidered accordingly.\par
There is strong counterevidence against interpreting the first example as a nonexistence result. On a global hyperbolic lattice, the integral isometry \(\operatorname{diag}(A,A^{-t})\), with \(A\) the unimodular companion matrix of \(f\), has characteristic polynomial \(f^2\). Adjoining \(-1\) on \(H^2\) realizes \(F\) globally. The demonstrated failure is the proposed obstruction construction, not nonexistence for this input, and not a proved failure of the Salem-realization theorems.\par
\section*{An excessive determinant quantifier at \(2\)}
Problem 2 (Minor).  Claim 2 in the proof of Proposition 2.6 (p. 9) allows arbitrary units \(u_+,u_-\) with \(u_+u_-=(-1)^n\), even when only one linear eigenspace exists. The determinant conclusion in that case is false.\par
Take \[F=(X-1)^2,\quad n=1,\quad n_+=2,\quad n_-=0,\quad
 F_1=1,\quad u_+=1,\quad u_-=-1.\] All stated numerical premises hold. The claim would supply an even unimodular rank-two \(\mathbb Z_2\)-lattice of determinant squareclass \(1\). Its Gram matrix would have the form \[\begin{pmatrix}2a&b\\b&2c\end{pmatrix},\qquad b\text{ odd}.\] Its determinant is \(4ac-b^2\equiv3\pmod4\), whereas a unit in squareclass \(1\) is \(1\bmod8\). This is impossible.\par
The same example disproves the sentence on p. 11 asserting that a nonzero linear eigenspace forces \(v_2(F_1(1))+v_2(F_1(-1))\) to be odd: here this sum is zero. The single-eigenspace choice on p. 9 also uses the exponent \(n_1\) where the asserted total determinant requires \(n\).\par
The ordinary existence conclusion survives this test: \(H\) with the identity is an even unimodular realization of \((X-1)^2\). In a single-eigenspace version, remove the arbitrary determinant prescription and choose the determinant compatible with the desired total form. Proposition 11.1 (p. 28) prescribes both component determinants only when both eigenspaces occur, so this example is not a counterexample to that proposition. The cited \(2\)-adic lattice criterion and the corrected single-eigenspace proof still require an explicit reconciliation.\par
\section*{The abstract signature characterization}
Problem 3 (Minor).  Proposition 3.7 (p. 13) overstates uniqueness in (a), and its listed conditions omit the hyperbolic constraint for nonreal reciprocal pairs off the unit circle.\par
For the identity on positive definite \(\mathbb R^4\), its signature map is \((4,0)\) on a polynomial vanishing at \(1\) and zero otherwise. Both \[(X-1)^4\quad\text{and}\quad(X-1)^2(X+1)^2\] are monic symmetric degree-four polynomials satisfying exactly the conditions stated in (a). Thus (a) alone does not give a unique polynomial, even with constant term \(1\). The multiplicity condition (d), applied jointly, selects the first polynomial.\par
There is also a distinct sufficiency issue. Let \[P=(X^2+4)(X^2+\tfrac14),\qquad
 \tau(u)=
 \begin{cases}(3,1)&P\mid u,\\(0,0)&P\nmid u.\end{cases}\] On nonzero real polynomials this satisfies the printed algebraic conditions (a)–(e). Condition (a) selects \(P\) uniquely. In (b), coprime symmetric factors cannot split the reciprocal pair \(X^2+4,\ X^2+\tfrac14\), so additivity holds. Taking the symmetric radical preserves divisibility by \(P\). Every irreducible symmetric real polynomial and every real reciprocal linear pair appearing in (d)–(e) has value zero and multiplicity zero.\par
However, an isometry annihilated by \(P\) on a nondegenerate four-dimensional space has the two reciprocal two-dimensional primary spaces totally isotropic. The form therefore has signature \((2,2)\), not \((3,1)\). The assertion is false even after the natural treatment of the zero polynomial; the example does not rely on the additional undefined expression \(\operatorname{rad}(0)\).\par
Move uniqueness after the joint multiplicity requirements, and add the hyperbolic multiplicity rule for all non-self-reciprocal factors, including nonreal quadratics. A complete reformulation should also specify the domain at the zero polynomial. The Salem-plus-cyclotomic applications have no off-circle nonreal quartet, so this example does not refute those applications.\par
\section*{Two dimension inconsistencies}
Problem 4 (Minor).  Two explicit constructions have dimensions incompatible with their claimed K3 signatures.\par
In the proof of Theorem 19.1, Case 1 (p. 42), for \(m=5\) or \(8\) the displayed choice is \[C=\Phi_m^2(X-1)^2.\] Since \(\varphi(5)=\varphi(8)=4\), this has degree \(10\), not the required \(12\). With \(\deg S=10\), the displayed \(SC\) has degree \(20\) and cannot be the polynomial of a signature map of maximum \((3,19)\). Replacing the final exponent by \(4\) repairs the degree and (C1): \(C(1)=0\), while \(C(-1)\) is respectively \(16\) or \(64\). The edge between \(S\) and \(\Phi_m\) is assumed; the usual ramified prime of \(\Phi_5\) or \(\Phi_8\) links it to \(X-1\) with multiplicity four. This checks the numerical correction, not every subsequent gluing assertion.\par
In the proof of Proposition 27.2 (p. 57), the complementary lattice is assigned signature \((1,22-\deg C)\). Its rank is then \(23-\deg C\), and its direct sum with the stated \((2,\deg C-2)\) lattice has signature \((3,20)\) and rank \(23\). In the branch \(\deg C\equiv4\bmod8\), the displayed complement does not even satisfy the necessary signature congruence. The required signature is \[(1,21-\deg C),\] giving rank \(22-\deg C\) and total signature \((3,19)\). This repair has the required congruence for \(\deg C\equiv4\bmod8\). Neither arithmetic error proves nonexistence of the intended automorphisms.\par
\section*{A missing root hypothesis}
Problem 5 (Minor).  Theorem 21.3 (p. 46) assumes only \(\delta\in\mathbb C^\times\) and \(|\delta|=1\), omitting \(S(\delta)=0\).\par
Choose any degree-four Salem polynomial, for example \[S=X^4-X^3-X^2-X+1.\] Its trace polynomial is \(Y^2-Y-3\), with one root \(>2\) and one in \((-2,2)\); irreducibility follows since the norm of \(r^2-4\) is \(-3\), not a square. Thus it satisfies the degree condition in part (a). Let \(\delta\) be a primitive \(29\)th root of unity. The printed hypotheses hold, but \(\delta\) cannot be an eigenvalue of an integral automorphism on the rank-\(22\) lattice \(H^2(X,\mathbb Z)\): its minimal polynomial has degree \(\varphi(29)=28\).\par
The introduction’s Theorem 3 (p. 2) and Notation 15.3 (p. 35) already contain the intended root restriction. Restoring \(S(\delta)=0\) aligns the local statement with them. This is a counterexample to the printed unrestricted quantifier, not to the introduction’s correctly restricted assertion.\par
\section*{A fixed characteristic polynomial obstructed by the root chamber}
Problem 6 (Major).  Theorem 24.1(a) (pp. 51–52), with the fixed polynomial \(F=S(X^2-1)\) specified immediately before it, is false. For the Salem polynomial below the discriminant action forced by the ambient unimodular lattice is incompatible with type (P).\par
Take \[S(X)=X^{20}+X^{19}-\sum_{j=3}^{17}X^j+X+1.\] Then \(S(1)=-11\) and \(S(-1)=1\), so the oddness, square-free and coprimality hypotheses hold. Here, as elsewhere in the theorem, square-free integers include \(1\).\par
\subsection*{An exact certificate that \(S\) is Salem}
Writing \(S(X)=X^{10}R(X+X^{-1})\) gives \[\begin{split}
R(Y)={}&Y^{10}+Y^9-10Y^8-10Y^7+34Y^6+33Y^5\\
 &-45Y^4-40Y^3+19Y^2+13Y-1.
\end{split}\] Exact integer Horner-equivalent evaluations, with denominators cleared by \(1000^{10}\), give a strict sign change on each of the following disjoint intervals; all endpoints shown are to be divided by \(1000\): \[\begin{array}{ccccc}
[-1973,-1972]&[-1777,-1776]&[-1519,-1518]&[-1170,-1169]&[-603,-602]\\
{}[70,71]&[750,751]&[1348,1349]&[1792,1793]&[2077,2078].
\end{array}\] There are ten intervals for a polynomial of degree ten, so they exhaust its roots. Nine roots lie in \((-2,2)\) and the remaining root exceeds \(2\).\par
For irreducibility, represent a polynomial over \(\mathbb F_2\) by the hexadecimal integer whose bit \(j\) is its \(X^j\) coefficient. The reduction of \(S\) is `1bfffb`. Exact modular squaring and Euclidean division give \[X^{2^{20}}\equiv X\pmod{\bar S},\qquad
 X^{2^{10}}\bmod\bar S=\mathtt{dfffd},\] and \[\gcd(\bar S,X^{2^{10}}-X)=
 \gcd(\bar S,X^{16}-X)=1.\] The prime divisors of \(20\) are \(2\) and \(5\), so these are the full irreducibility test over \(\mathbb F_2\). Thus \(S\) is irreducible over \(\mathbb Q\). Together with the trace-root intervals this proves that \(S\) is a Salem polynomial of degree \(20\). The complete integer endpoint values and modular remainders, independently recomputed by the primary reviewer using exact BigInt arithmetic, are saved in `exact\_salem\_check.json`. No floating-point approximation or author program enters this certificate.\par
\subsection*{The discriminant and positivity contradiction}
Suppose that the asserted realizable isometry \(t\) exists. Its cyclotomic lattice \[N=L\cap\ker(t^2-1)\] is even, negative definite and primitive of rank two. The orthogonal projection of the unimodular lattice \(L\) onto \(N\otimes\mathbb Q\) is \(N^*\). Since \(S(t)\) kills the Salem subspace, \[S(t)N^*\subset N.\] The determinant of \(S(t)\) on \(N\otimes\mathbb Q\) is \(S(1)S(-1)=-11\). Comparing indices therefore gives \[|\det N|\mid11.\] There is no even unimodular rank-two definite lattice, hence \(|\det N|=11\).\par
The positive lattice \(N(-1)\) is uniquely represented by \[B=\begin{pmatrix}2&1\\1&6\end{pmatrix}.\] Indeed, write a reduced positive even binary form as \(\left(\begin{smallmatrix}2a&b\\b&2c\end{smallmatrix}\right)\) with \(|b|\le a\le c\). Then \(11=4ac-b^2\ge3a^2\), so \(a=1\), \(b=\pm1\) and \(c=3\). Its only roots are \(\pm e_1\), since \[x^2+xy+3y^2=(x+y/2)^2+\tfrac{11}{4}y^2=1\] forces \(y=0\). Type (P), together with characteristic polynomial \(X^2-1\), forces \(t(e_1)=e_1\) and hence \[t=\begin{pmatrix}1&1\\0&-1\end{pmatrix}.\] On the discriminant group of order \(11\), generated by \((-e_1+2e_2)/11\), this acts by \(-1\).\par
On the other hand, reducing \(S\) modulo \(X^2-1\) gives \[S(t)=-5-6t.\] The containment \(S(t)N^*\subset N\) says that this expression acts as zero on the discriminant group. Modulo \(11\), it forces \(t=+1\), contradicting the required \(-1\) action. No realization with the theorem’s fixed polynomial and Salem signature exists.\par
The failure is at the final positivity assertion in the proof on p. 52, not merely at an unproved numerical irreducibility assertion. A repair must restrict the fixed-polynomial assertion, add the needed discriminant/chamber compatibility, or allow the cyclotomic polynomial to change. This counterexample does not disprove the general realizability of the Salem number after such a change, including the broader degree-\(20\) result attributed to Takada. The Weyl correction discussed below explains why that distinction matters.\par
\section*{Positive checks, counterevidence and remaining limits}
The following checks help delimit the findings.\par
First, the real Hasse comparison in Proposition 9.3 depends on the negative index modulo four, since \(w_2(\langle1^r,-1^s\rangle)\) is determined by \(\binom{s}{2}\bmod2\). This supports the stated invariance when the positive complex place of a fixed Salem factor is moved. It does not repair the finite-place obstruction in Problem 1.\par
Second, an apparent gap in passing from equal Witt residues to stability of a particular glued lattice is not by itself a counterexample to chamber realization. If an integral isometry \(t\) of \(L\) has negative definite cyclotomic part \(L_C\), each root \(r\in L_C\) gives the integral reflection \[s_r(x)=x+(x,r)r.\] It fixes the rational Salem part pointwise. The finite Weyl group acts transitively on the chambers of \(L_C\); choose \(w\) so that \(wt\) preserves a chosen chamber. Then \(wt\) still preserves \(L\), agrees with \(t\) on the Salem part, and has finite-order restriction to \(L_C\). It is semisimple and has the same Salem signature. This supplies a concrete existence mechanism when the cyclotomic polynomial may change, as allowed in Definition 15.5. A direct appeal to Proposition 16.1 for a product of different cyclotomic factors on p. 42 exceeds that proposition’s single-power premise; the reflection argument explains a way to avoid treating this compressed step as a proved nonexistence defect.\par
Third, Theorem 21.4 uses the positive period plane and the chamber condition. For a simple unit-circle root of the irreducible Salem factor, a rational vector in the Salem space orthogonal to that plane would generate, under \(\mathbb Q[t]\), the whole Salem space; that is incompatible with orthogonality to the nonzero plane. Hence the rational Picard part is the cyclotomic part as required. Positive entropy also makes the remaining orthogonal part negative definite in the nonprojective case, supporting semisimplicity in Proposition 20.6.\par
Fourth, the argument of Proposition 25.5 correctly uses the simple \(\delta\)-eigenline to recover \(H^{2,0}\) from the isometric structure. Weak Torelli identifies the surfaces; it need not identify the original chosen isometry with a cone-preserving automorphism. Finiteness still imports the integral isometric-structure theorem cited as [BM 79]. The finite cyclotomic possibilities follow from the fixed degree bound. The restrictions in Section 26 and the resultant-one splitting in Lemma 27.1 were traced, but numerical tables, complete root-system classifications and all cited global theorems were not independently reproduced.\par
No author program was executed. The degree-twenty certificate above was computed with an independently written exact-arithmetic calculation. Bounded searches of the publisher, arXiv and title-specific code/correction results found no article-specific supplementary code or erratum; unrelated K3 software is not treated as this paper’s code. Absence of code is not a defect. This review does not certify a corrected global obstruction theorem, every arithmetic example, or equivalence with the publisher text.\par
\section*{Supporting checks}
Three native helpers supplied bounded text checks: Part I, signature/surgery, and K3 realization/finiteness. Their precise actual reading ranges and limitations are preserved in the accompanying record; none is represented as an independent acceptance or as a verified separate scientific model. The controller’s UI observation is recorded separately from the runtime’s self-description. Source, TeX, compiler, dependencies, logs, PDF and image bytes and hashes, actual exit codes, and final report-page inspection are bound by the build and delivery records. The workflow state remains `scientific\_decision=pending` with `human\_reviews=[]`.\par
Conclusion.\par
Item | State\par
Supported findings | Two Major and four Minor findings. Obstruction choice independence and the fixed degree-20 realization assertion fail as stated.\par
Repair verification | Exact counterexamples and bounded local corrections checked; replacement obstruction theory and complete revised gluing/classification remain uncertified.\par
Author adoption | Not known; no communication sent.\par
Remaining work | Rebuild the obstruction argument, correct fixed-polynomial chamber compatibility, and verify downstream arguments.\par
Scientific decision | Pending; human reviews empty.\par
9 E. Bayer-Fluckiger, Automorphisms of K3 surfaces, signatures, and isometries of lattices, arXiv:2209.06698v4, 1 May 2024, 59 pp. https://arxiv.org/abs/2209.06698v4. EMS Press, article metadata for DOI 10.4171/JEMS/1598, online first 30 January 2025. https://ems.press/journals/jems/articles/14298545. Metadata only; publisher text not assumed equivalent.\par
\end{document}
